% AI-written, AI cross-reviewed; no human sign-off; not a claim that RH is proved
% Manuscript v3, heading p. 148; original paper2/pending_p096.tex, SHA-256 afa227cf428f335f3d9be5e09c8e3f51c12c3d058ff5cbc6ebbd1f0445574726.
\subsection{Candidate proof of equality away from the moved place (P096)}
\label{sec:pending-p096}
\noindent\textsf{P096. The Round 21 AI reviewer found the statement and argument correct and referred the candidate for separate cross-review.}
Fix the data of \nodeid{P096}, put $n=N_i$, and write $v_0$ for the
place of $F$ specified by the move. Let $L_c,L_{c'}\subset E^n$ be
the two class lattices in the chosen coordinates.

\begin{proof}[Candidate proof]
If $i=0$ or $i=m$, the first clause of \code{AdC} supplies the
endpoint level $U$ and representatives $y,y'$ of $c,c'$ in that
level. If $0<i<m$, the second clause gives
$\operatorname{chOf}_i(c')=\operatorname{chOf}_i(c)$. Use the
chart at this common label, its level $U$, and its two representatives.
In either case both classes use the same base lattice $B$. The level
field \nolinkurl{Lb_fin} gives finite generation over $\mathcal O_E$,
so its generators also generate the integral span after local extension.
Writing a hat for the span over the integral finite adelic coefficient ring,
the \nolinkurl{y_lat} field of the endpoint level or of the common chart gives
$\widehat L_c=y\widehat B$ and $\widehat L_{c'}=y'\widehat B$.
The one-place relation gives
\[
    \gamma_f y'=y\eta\kappa,\qquad
    \kappa\in U,\qquad \eta_w=1\quad(w\ne v_0),
\]
where $\gamma_f$ is the diagonal image of the rational matrix.
At each $w$ the local level lies in the stabilizer of the integral
local span of $B$. Thus $\kappa_w$ preserves that span in both
directions, since $\kappa_w^{-1}$ also belongs to the local level.

Fix $Q$ over $w\ne v_0$. For a rational lattice $L$ write
$L_Q^\wedge=\mathcal O_{E_Q}\cdot L\subset E_Q^n$. Evaluating adelic
coordinates at $w$ gives coefficients in $E\otimes_F F_w\simeq\prod_{Q\mid w}E_Q$.
Projection to the factor indexed by $Q$ sends the integral coefficients
into $\mathcal O_{E_Q}$: the generators from $\mathcal O_E$ and from
the integral adeles of $F$ are integral there. The
$\mathcal O_{E_Q}$-span of the evaluated lattice is exactly
$L_Q^\wedge$. Indeed its generators include the image of $L$, and
every evaluated linear combination belongs to $\mathcal O_{E_Q}\cdot L$.
The same reasoning applies to the integral local span at $w$. The
representative identities therefore give
$L_{c,Q}^\wedge=y_QB_Q^\wedge$ and
$L_{c',Q}^\wedge=y'_QB_Q^\wedge$, with
$\kappa_Q B_Q^\wedge=B_Q^\wedge$. Evaluation of the displayed
one-place relation now yields
\[
  (\gamma L_{c'})_Q^\wedge
     =\gamma y'_Q B_Q^\wedge
     =y_Q\kappa_Q B_Q^\wedge
     =L_{c,Q}^\wedge .
\]

To descend the equality, put $R=(\mathcal O_E)_Q$. For any full finite
lattice $L\subset E^n$, $L_Q=R L$ is finite and torsion free over the
DVR $R$. An $R$-basis is therefore also an $E$-basis of $E^n$.
In that basis a rational vector in $L_Q^\wedge$ has coefficients in
$E\cap \mathcal O_{E_Q}=R$, since the completed valuation restricts to
the $Q$-valuation on $E$. Consequently $E^n\cap L_Q^\wedge=L_Q$.
Both class lattices are finite and span $E^n$ by the class-lattice
hypothesis and the coordinate isomorphism. The invertible rational
matrix $\gamma$ preserves these two properties. Intersecting the
completed equality with $E^n$ gives $(\gamma L_{c'})_Q=(L_c)_Q$.
\end{proof}

\noindent\small The two clauses are in
\nolinkurl{IwasawaDecomposition/MemSmulLk.lean:397--403} and the adelic
identity in \nolinkurl{HermitianColumns/CoordinateLatticeSelfDual.lean:17--21}.
The general argument and its specialization \code{use₆} are in
\nolinkurl{IwasawaDecomposition/AdelicConjugateCharacterStandardForms.lean:206--241}.
The local product identification is \nolinkurl{GroundTheory.Global.SemiLocal.tensorEquiv'}
in \nolinkurl{GroundTheory/GlobalFields/SemilocalAlgebras/Finite.lean:111--119}.
\normalsize
